\begin{textsample}{NEG Dim 5 – vem\_ed\_17 – Score -2.00 – vem\_ed\_17/t080.txt}  \label{ex:f5_neg_031}
BOAS PRÁTICAS 2) Fatec Bragança Paulista e Uniminuto (Colômbia) – PCI “Lengua Española” envolve alunos de Gestão Empresarial e a turma de Licenciatura em Língua Espanhola conduzida por Margot Costas. “As tarefas foram pensadas a partir do perfil dos estudantes de ambas as instituições.
A primeira atividade é individual, a tarefa principal reúne os dois países e a última fica a cargo dos alunos colombianos, que ministrarão pequenas aulas de língua espanhola para os brasileiros.
Essa é uma proposta inovadora, que nunca fizemos em edições anteriores, por isso estamos otimistas com o resultado”, empolga-se Lilian.
Site do PCI: sites.google.com/view/eciaespanol/lengua-espa\%C3\%B1ola 3) Fatec Bragança e Universidad Distrital Francisco José de Caldas (Colômbia) – PCI “Entreculturas” conta com alunos brasileiros de Logística e colombianos que estudam Educação Infantil com o professor Miguel Nicholls.
Após as apresentações individuais em vídeo, há duas tarefas em equipes mistas, enfocando aspectos culturais de ambos os países. “Os alunos brasileiros praticam a língua espanhola de forma efetiva com os falantes nativos”, ressalta.
Site do PCI: sites.google.com/view/eciadistrital/entreculturas 4) Fatec Itu com DUOC (Chile) – PCI “Aspectos Sociales” une estudantes de Gestão Empresarial e de \textbf{Avaliação} de Projetos (professor Felipe Zambrano).
As tarefas ocorrem ao longo de um ano, algo inovador no planejamento de PCIs: no primeiro semestre, serão solicitadas atividades de rompehielo (quebra-gelo), tarefa intercultural e um levantamento das problemáticas sociais nas regiões em que os alunos vivem.
Isso prepara o cenário para a tarefa principal a ser elaborada no segundo semestre: um aplicativo sobre essas problemáticas.
Site do PCI: sites.google.com/view/eciachile/aspectos-sociales Sobre o aprendizado principal com esses projetos, Lilian sintetiza os aspectos mais relevantes: “O enfoque maior é a interação com os estrangeiros, de forma a propiciar aos alunos essa vivência além-fronteiras, usando a língua estudada em sala de aula.
Os aspectos culturais também são determinantes.
Percebeu-se, nesses anos de PCI, quanto os alunos se desenvolvem linguística e culturalmente com essa experiência”.

% matched lemmas: avaliação
\end{textsample}
